% dynamics/state_taxonomy.tex

\chapter{State Taxonomy}

\label{chap:state-taxonomy}

% ============================================================
\section{Motivation}
% ============================================================

\begin{remark}
Within AIM, the term \emph{state} appears in multiple contexts:
\begin{itemize}
\item geometric reconstruction,
\item runtime evaluation,
\item realization,
\item operational dynamics,
\item and vehicle interpretation.
\end{itemize}
\end{remark}

\begin{remark}
Without a consistent taxonomy, different state concepts may become
ambiguous or unintentionally conflated.
\end{remark}

\begin{remark}
This chapter therefore introduces a structural classification of state
concepts within AIM.
\end{remark}

% ============================================================
\section{General State Interpretation}
% ============================================================

\begin{definition}[State]
\label{def:taxonomy-general-state}
A state is a structured description of a system configuration relative
to a specified interpretation layer.
\end{definition}

\begin{principle}[State-layer principle]
\label{princ:state-layer-principle}
Within AIM, state concepts belong to distinct interpretation layers and
must not be identified implicitly.
\end{principle}

% ============================================================
\section{Intrinsic Geometric States}
% ============================================================

\subsection{pose2}

\begin{definition}[pose2 state]
\label{def:taxonomy-pose2}
The pose2 state is the intrinsic planar reconstruction state generated
from:
\[
\kappa(s).
\]
\end{definition}

\begin{remark}
pose2 represents:
\begin{itemize}
\item intrinsic geometry,
\item planar tangent evolution,
\item and reconstruction state.
\end{itemize}
\end{remark}

\begin{remark}
pose2 is fundamentally an intrinsic geometric state.
\end{remark}

% ------------------------------------------------------------

\subsection{pose3}

\begin{definition}[pose3 state]
\label{def:taxonomy-pose3}
The pose3 state is the runtime-evaluated spatial railway state generated
from:
\begin{itemize}
\item pose2 reconstruction,
\item profile evaluation,
\item cant interpretation,
\item frame evaluation,
\item and runtime coordinate interpretation.
\end{itemize}
\end{definition}

\begin{remark}
pose3 describes the railway runtime state rather than the vehicle
itself.
\end{remark}

\begin{remark}
pose3 therefore belongs to the runtime railway-state layer.
\end{remark}

% ============================================================
\section{Runtime States}
% ============================================================

\begin{definition}[Runtime state]
\label{def:taxonomy-runtime-state}
A runtime state is a dynamically evaluated state generated during
runtime interpretation.
\end{definition}

\begin{remark}
Runtime states depend on:
\begin{itemize}
\item runtime operators,
\item realization,
\item coordinate interpretation,
\item and operational context.
\end{itemize}
\end{remark}

\begin{remark}
Examples include:
\begin{itemize}
\item pose3,
\item runtime frames,
\item projection states,
\item and operational interpretation states.
\end{itemize}
\end{remark}

% ============================================================
\section{Realized States}
% ============================================================

\begin{definition}[Realized state]
\label{def:taxonomy-realized-state}
A realized state is a runtime state generated from physically realized
rather than purely analytical geometry.
\end{definition}

\begin{remark}
Realized states depend on:
\begin{itemize}
\item construction,
\item measurement,
\item maintenance,
\item realization filtering,
\item and operational degradation.
\end{itemize}
\end{remark}

\begin{remark}
Typical examples include:
\[
\mathrm{pose3}_{\mathrm{real}},
\qquad
\kappa_{\mathrm{real}},
\qquad
\phi_{\mathrm{real}}.
\]
\end{remark}

% ============================================================
\section{Operational States}
% ============================================================

\begin{definition}[Operational state]
\label{def:taxonomy-operational-state}
An operational state is a runtime state relevant for railway operation,
comfort, admissibility, or dynamic interpretation.
\end{definition}

\begin{remark}
Operational states may include:
\begin{itemize}
\item effective acceleration,
\item jerk,
\item roll evolution,
\item cant deficiency,
\item and balance trajectories.
\end{itemize}
\end{remark}

\begin{remark}
Operational states are derived from runtime railway states rather than
stored directly as intrinsic geometry.
\end{remark}

% ============================================================
\section{Vehicle States}
% ============================================================

\begin{definition}[Vehicle state]
\label{def:taxonomy-vehicle-state}
A vehicle state is a runtime-relative operational state describing the
vehicle relative to the railway runtime state.
\end{definition}

\begin{remark}
Vehicle states depend on:
\begin{itemize}
\item railway runtime state,
\item operational velocity,
\item vehicle interpretation,
\item and dynamic response.
\end{itemize}
\end{remark}

\begin{remark}
Thus:
\[
\mathrm{vehicleState}
=
\mathcal V(
\mathrm{pose3},
v,
\dots
).
\]
\end{remark}

\begin{remark}
Vehicle states are therefore not intrinsic railway states.
\end{remark}

% ============================================================
\section{Center-of-Gravity States}
% ============================================================

\begin{definition}[Center-of-gravity state]
\label{def:taxonomy-cg-state}
A center-of-gravity state is an operational abstraction describing the
runtime trajectory of the vehicle center of gravity relative to the
railway runtime state.
\end{definition}

\begin{remark}
Center-of-gravity states are primarily operational interpretation states
rather than full mechanical vehicle models.
\end{remark}

\begin{remark}
These states are particularly relevant for:
\begin{itemize}
\item comfort interpretation,
\item balance interpretation,
\item Schwerpunkt-Trassierung,
\item and runtime operational optimization.
\end{itemize}
\end{remark}

% ============================================================
\section{Dynamic States}
% ============================================================

\begin{definition}[Dynamic state]
\label{def:taxonomy-dynamic-state}
A dynamic state is a runtime state participating in operational state
evolution.
\end{definition}

\begin{remark}
Dynamic states evolve according to:
\[
x'(s)
=
f(x,u).
\]
\end{remark}

\begin{remark}
Dynamic states may include:
\begin{itemize}
\item acceleration,
\item jerk,
\item roll evolution,
\item and operational balance quantities.
\end{itemize}
\end{remark}

% ============================================================
\section{Relative States}
% ============================================================

\begin{definition}[Relative state]
\label{def:taxonomy-relative-state}
A relative state is a state evaluated relative to another runtime state
or runtime reference frame.
\end{definition}

\begin{remark}
Within AIM, many operational quantities are fundamentally relative
states.
\end{remark}

\begin{remark}
Examples include:
\begin{itemize}
\item vehicle states relative to track states,
\item wheelset states relative to local frames,
\item and center-of-gravity trajectories relative to realized track
geometry.
\end{itemize}
\end{remark}

% ============================================================
\section{State Transformations}
% ============================================================

\begin{remark}
Different state classes are connected through runtime operators.
\end{remark}

\begin{remark}
Examples include:
\[
\mathcal R,
\qquad
\mathcal D,
\qquad
\mathcal V.
\]
\end{remark}

\begin{remark}
Typical transformations include:
\begin{itemize}
\item intrinsic reconstruction,
\item realization filtering,
\item runtime evaluation,
\item vehicle interpretation,
\item and operational state evolution.
\end{itemize}
\end{remark}

% ============================================================
\section{State Hierarchy}
% ============================================================

\begin{remark}
The AIM state hierarchy may be interpreted schematically as:
\[
\kappa(s)
\rightarrow
\mathrm{pose2}
\rightarrow
\mathrm{pose3}
\rightarrow
x(s)
\rightarrow
x_{\mathrm{vehicle}}
\rightarrow
x_{\mathrm{CG}}.
\]
\end{remark}

\begin{remark}
Realization operators act throughout this hierarchy:
\[
\mathcal R :
x
\mapsto
x_{\mathrm{real}}.
\]
\end{remark}

% ============================================================
\section{Canonical Interpretation}
% ============================================================

\begin{proposition}
Within AIM, state concepts belong to distinct interpretation layers.
\end{proposition}

\begin{proposition}
pose3 describes the railway runtime state rather than the vehicle
itself.
\end{proposition}

\begin{proposition}
Vehicle states are relative runtime interpretation states generated from
interaction with the railway runtime state.
\end{proposition}

\begin{proposition}
Operational and dynamic states are generated through runtime evaluation
rather than through static geometric storage.
\end{proposition}

% ============================================================
\section{Connection to AIM}
% ============================================================

\begin{summary}
State taxonomy provides the structural interpretation framework for AIM
runtime architecture.

Within this framework:
\begin{itemize}
\item pose2 represents intrinsic reconstruction state,
\item pose3 represents railway runtime state,
\item realized states represent physical railway realization,
\item operational states describe runtime behaviour,
\item vehicle states describe relative operational interpretation,
\item and dynamic states describe runtime state evolution.
\end{itemize}

This taxonomy prevents conceptual ambiguity between:
\begin{itemize}
\item geometry,
\item runtime,
\item realization,
\item dynamics,
\item and vehicle interpretation.
\end{itemize}
\end{summary}
